\documentclass[11pt]{article}
\usepackage{amsmath, amssymb, amsthm}
\usepackage{geometry}
\geometry{margin=1in}

\title{T-FAN Supplementary Material}
\author{Validation and Proofs Team}
\date{Updated Section D -- November 2025}

\begin{document}
\maketitle

\section{Notation Recap}
We denote by $\mathcal{F}$ the family of admissible fan profiles, parameterized by the curvature radius $r$ and local twist $\theta$. The aerodynamic response under turbulent forcing $\boldsymbol{u}(t)$ is represented by the linear operator $\mathcal{T}$ acting on the perturbation field $\phi$.

\section{Baseline Stability Bound}
Given the energy functional
\begin{equation}
    E(\phi) = \frac{1}{2} \int_{\Omega} \left( \|\nabla \phi\|^2 + \alpha \phi^2 \right) \mathrm{d}x,
\end{equation}
we recall the bound derived in the main text:
\begin{equation}
    E(\phi(t)) \leq E(\phi(0)) e^{-2\lambda t} + \frac{\|\boldsymbol{u}\|^2}{2\lambda} \left( 1 - e^{-2\lambda t} \right),
\end{equation}
where $\lambda$ is the spectral gap of the generator associated with $\mathcal{T}$.

\section{Proof of Proposition C.3}
The proof follows from compactness arguments and the coercivity of the bilinear form $B(\phi, \psi)$. Specifically, let $\{\phi_n\}$ be the minimizing sequence. The reflexivity of $H^1(\Omega)$ implies the existence of a weakly convergent subsequence $\phi_{n_k} \rightharpoonup \phi_*$, and the lower semi-continuity of $E$ ensures that $\phi_*$ is a minimizer.

\section{Updated Section D: Turbulence-Adjusted Fan Stability}
\label{sec:updated-section-d}
We refine the previous analysis by incorporating the turbulence intensity parameter $\tau$. Define the modified operator
\begin{equation}
    \mathcal{T}_\tau = \mathcal{T} + \tau \mathcal{K},
\end{equation}
where $\mathcal{K}$ is compact and self-adjoint, representing the feedback from the turbulence control surface.

\subsection{Spectral Adjustment}
Let $\{\mu_i\}$ denote the eigenvalues of $\mathcal{T}$ ordered non-decreasingly. Weyl's inequality for compact perturbations yields
\begin{equation}
    |\mu_i(\mathcal{T}_\tau) - \mu_i(\mathcal{T})| \leq \tau \|\mathcal{K}\|.
\end{equation}
Consequently, the perturbed spectral gap satisfies
\begin{equation}
    \lambda_\tau \geq \lambda - \tau \|\mathcal{K}\|,
\end{equation}
which is positive provided $0 \leq \tau < \lambda / \|\mathcal{K}\|$.

\subsection{Energy Decay Under Turbulence}
Under the modified dynamics, the energy functional obeys
\begin{equation}
    \frac{\mathrm{d}}{\mathrm{d}t} E(\phi) \leq -2\lambda_\tau E(\phi) + \|\boldsymbol{u}\| \cdot \|\phi\|.
\end{equation}
Applying Grönwall's inequality and the bound on $\lambda_\tau$ leads to
\begin{equation}
    E(\phi(t)) \leq E(\phi(0)) e^{-2(\lambda - \tau \|\mathcal{K}\|) t} + \frac{\|\boldsymbol{u}\|^2}{2(\lambda - \tau \|\mathcal{K}\|)} \left( 1 - e^{-2(\lambda - \tau \|\mathcal{K}\|) t} \right).
\end{equation}
This explicit decay rate quantifies the resilience of T-FAN to turbulent fluctuations of magnitude $\tau$.

\subsection{Critical Intensity Threshold}
Finally, observe that if $\tau$ approaches the critical value $\lambda/\|\mathcal{K}\|$, the decay rate vanishes. In that regime, higher-order stabilizers from the nonlinear actuator terms must be included. Empirical results (see Dataset TF-2025A) confirm that operational turbulence remains below $0.6\,\lambda/\|\mathcal{K}\|$, validating the linear theory.

\section{Conclusion}
The updated Section~\ref{sec:updated-section-d} establishes explicit quantitative bounds that align with the November 2025 experimental campaign. The synergy between theoretical guarantees and empirical validation strengthens the certification path for T-FAN.

\end{document}
